\subsection{Moments, Sample Statistics, and Random Vectors}
\begin{definition} \label{def:moments-correlation-and-sample-statistics} \hlt{Moments, Correlation, and Sample Statistics}
\begin{enumerate}[label=(\roman*)]
\setlength{\itemsep}{0pt}
\item The \hlt{$k$th moment} of $X$ is $\mathbb{E}(X^k)$, assuming $\mathbb{E}|X|^k<\infty$.
\item The \hlt{correlation} of $X,Y$ is $\operatorname{Corr}(X,Y)\defeq\frac{\operatorname{Cov}(X,Y)}{\sqrt{\operatorname{Var}(X)}\sqrt{\operatorname{Var}(Y)}}$, assuming $\operatorname{Var}(X),\operatorname{Var}(Y)>0$.
\item The \hlt{sample mean} of $X_1,\ldots,X_n$ is $\overline{X}_n\defeq n^{-1}\sum\limits_{i=1}^{n}X_i$.
\item The \hlt{sample variance} of $X_1,\ldots,X_n$ is $S_n^2\defeq(n-1)^{-1}\sum\limits_{i=1}^{n}(X_i-\overline{X}_n)^2$.
\end{enumerate}
\end{definition}

\begin{theorem} \label{thm:expectation-variance-covariance-rules} \hlt{Expectation, Variance, and Covariance Algebra}\\
Assume all displayed quantities exist and $a_i,a,b\in\R$.
\begin{enumerate}[label=(\roman*)]
\setlength{\itemsep}{0pt}
\item \hlt{Linearity}: $\mathbb{E}(\sum\limits_i a_iX_i)=\sum\limits_i a_i\mathbb{E}X_i$.
\item \hlt{Independent product}: if $X_1,\ldots,X_n$ are independent, then $\mathbb{E}(\prod\limits_i X_i)=\prod\limits_i\mathbb{E}X_i$.
\item \hlt{Variance identities}: $\operatorname{Var}(X)=\mathbb{E}(X^2)-(\mathbb{E}X)^2$ and $\operatorname{Var}(aX+b)=a^2\operatorname{Var}(X)$.
\item \hlt{Covariance identity}: $\operatorname{Cov}(X,Y)=\mathbb{E}(XY)-\mathbb{E}X\,\mathbb{E}Y$.
\item \hlt{Variance of a linear combination}: $\operatorname{Var}(\sum\limits_i a_iX_i)=\sum\limits_i a_i^2\operatorname{Var}(X_i)+2\sum\limits_{i<j}a_ia_j\operatorname{Cov}(X_i,X_j)$.
\end{enumerate}
\begin{proof}
\begin{enumerate}[label=(\roman*)]
\setlength{\itemsep}{0pt}
\item $\mathbb{E}(\sum\limits_i a_iX_i)=\int_\Omega\sum\limits_i a_iX_i\,d\mathbb{P}=\sum\limits_i a_i\int_\Omega X_i\,d\mathbb{P}=\sum\limits_i a_i\mathbb{E}X_i$.
\item By independence, $\mathcal{L}(X_1,\ldots,X_n)=\prod\limits_i F_{X_i}$; $\mathbb{E}(\prod\limits_i X_i)=\int_{\R^n}\prod\limits_i x_i\,d\prod\limits_i F_{X_i}(x_i)=\prod\limits_i\int_\R x_i\,dF_{X_i}(x_i)=\prod\limits_i\mathbb{E}X_i$.
\item $\operatorname{Var}(X)=\mathbb{E}(X^2-2X\mathbb{E}X+(\mathbb{E}X)^2)=\mathbb{E}(X^2)-(\mathbb{E}X)^2$.\\
Also $\operatorname{Var}(aX+b)=\mathbb{E}[(aX+b-a\mathbb{E}X-b)^2]=a^2\operatorname{Var}(X)$.
\item $\operatorname{Cov}(X,Y)=\mathbb{E}(XY-X\mathbb{E}Y-Y\mathbb{E}X+\mathbb{E}X\mathbb{E}Y)=\mathbb{E}(XY)-\mathbb{E}X\,\mathbb{E}Y$.
\item With $S=\sum\limits_i a_iX_i$, (iv) gives\\
$\operatorname{Var}(S)=\operatorname{Cov}(S,S)=\sum\limits_{i,j}a_ia_j\operatorname{Cov}(X_i,X_j)=\sum\limits_i a_i^2\operatorname{Var}(X_i)+2\sum\limits_{i<j}a_ia_j\operatorname{Cov}(X_i,X_j)$.
\end{enumerate}
\hlt{Structural Commentary}: Expectation is linear; variance and covariance are quadratic and bilinear consequences.\\
\hlt{Key Insight}: Expand products, then use linearity.
\end{proof}
\end{theorem}

\begin{example} \label{ex:indicator-sums-binomial-random-walk-game} \hlt{Indicator Sums as a Method}
\begin{enumerate}[label=(\roman*)]
\setlength{\itemsep}{0pt}
\item If $X\sim\operatorname{Binomial}(n,p)$, write $X=\sum\limits_{i=1}^{n}I_i$ with independent $I_i\sim\operatorname{Bernoulli}(p)$.\\
Then $\mathbb{E}X=np$ and $\operatorname{Var}(X)=np(1-p)$.
\item Let $\xi_i\in\{-1,1\}$ with $\mathbb{P}(\xi_i=-1)=p$ and $\mathbb{P}(\xi_i=1)=1-p$, independently, and $X_n=\sum\limits_{i=1}^{n}\xi_i$.\\
Then $\mathbb{E}X_n=n(1-2p)$ and $\operatorname{Var}(X_n)=4np(1-p)$.
\item Let $F_n=c\prod\limits_{i=1}^{n}M_i$ with independent $M_i\in\{2,\frac{1}{2}\}$ and $\mathbb{P}(M_i=2)=\mathbb{P}(M_i=\frac{1}{2})=\frac{1}{2}$.\\
Then $\mathbb{E}F_n=c\prod\limits_{i=1}^{n}\mathbb{E}M_i=c(\frac{5}{4})^n$.
\end{enumerate}
\end{example}

\begin{proposition} \label{prop:variance-zero-iff-constant} \hlt{Zero Variance Means Almost Surely Constant}\\
For $X\in L^2$ and $m\defeq\mathbb{E}X$, $\operatorname{Var}(X)=0$ iff $X=m$ a.s.
\begin{proof}
\begin{enumerate}[label=\roman*.]
\setlength{\itemsep}{0pt}
\item[$(\Rightarrow)$] If $\operatorname{Var}(X)=0$, then $\mathbb{E}[(X-m)^2]=0$. Since $(X-m)^2\geq0$, $\mathbb{P}((X-m)^2>0)=0$, hence $X=m$ a.s.
\item[$(\Leftarrow)$] If $X=m$ a.s., then $(X-m)^2=0$ a.s., so $\operatorname{Var}(X)=\mathbb{E}[(X-m)^2]=0$.
\end{enumerate}
\hlt{Structural Commentary}: Variance is the mean squared distance from the mean.\\
\hlt{Key Insight}: Zero nonnegative expectation forces zero almost surely.
\end{proof}
\end{proposition}

\begin{proposition} \label{prop:tail-integral-formula-nonnegative} \hlt{Tail Integral Formula}\\
If $X\geq0$ has density $f_X$, distribution function $F_X$, and $\mathbb{E}X<\infty$,\\
then $\mathbb{E}X=\int_0^\infty\mathbb{P}(X>x)\,dx=\int_0^\infty(1-F_X(x))\,dx$.
\begin{proof}
For $b>0$, integration by parts gives $\int_0^bxf_X(x)\,dx=-b(1-F_X(b))+\int_0^b(1-F_X(x))\,dx$.\\
Since $\mathbb{E}X<\infty$, $b(1-F_X(b))\leq\int_b^\infty xf_X(x)\,dx\to0$.\\
Letting $b\to\infty$ yields $\mathbb{E}X=\int_0^\infty(1-F_X(x))\,dx=\int_0^\infty\mathbb{P}(X>x)\,dx$.\\
\hlt{Structural Commentary}: A nonnegative variable is integrated by summing its survival probabilities.\\
\hlt{Key Insight}: Expectation equals area under the tail.
\end{proof}
\end{proposition}

\begin{example} \label{ex:maximum-uniform-expectation} \hlt{Maximum of Uniform Samples}\\
Let $X_1,\ldots,X_n$ be independent $\operatorname{Uniform}(0,1)$ variables and $Y_n=\max_iX_i$.\\
Then $F_{Y_n}(y)=\mathbb{P}(X_1\leq y,\ldots,X_n\leq y)=y^n$ for $0\leq y\leq1$.\\
So $\mathbb{E}Y_n=\int_0^1(1-F_{Y_n}(y))\,dy=\int_0^1(1-y^n)\,dy=\frac{n}{n+1}$.
\end{example}

\begin{theorem} \label{thm:sample-mean-and-variance-moments} \hlt{Mean and Variance of the Basic Sample Statistics}\\
Let $X_1,\ldots,X_n$ be iid with $\mathbb{E}X_i=\mu$ and $\operatorname{Var}(X_i)=\sigma^2<\infty$.
\begin{enumerate}[label=(\roman*)]
\setlength{\itemsep}{0pt}
\item \hlt{Sample-mean expectation}: $\mathbb{E}\overline{X}_n=\mu$.
\item \hlt{Sample-mean variance}: $\operatorname{Var}(\overline{X}_n)=\frac{\sigma^2}{n}$.
\item \hlt{Unbiased sample variance}: $\mathbb{E}S_n^2=\sigma^2$.
\end{enumerate}
\begin{proof}
\begin{enumerate}[label=(\roman*)]
\setlength{\itemsep}{0pt}
\item $\mathbb{E}\overline{X}_n=n^{-1}\sum\limits_{i=1}^{n}\mathbb{E}X_i=n^{-1}n\mu=\mu$.
\item Independence gives $\operatorname{Var}(\overline{X}_n)=n^{-2}\sum\limits_{i=1}^{n}\operatorname{Var}(X_i)=n^{-2}n\sigma^2=\frac{\sigma^2}{n}$.
\item Since $\sum\limits_{i=1}^{n}(X_i-\overline{X}_n)^2=\sum\limits_{i=1}^{n}(X_i-\mu)^2-n(\overline{X}_n-\mu)^2$. Thus $\mathbb{E}[(n-1)S_n^2]=n\sigma^2-n(\frac{\sigma^2}{n})=(n-1)\sigma^2$.
\end{enumerate}
\hlt{Structural Commentary}: Averaging preserves mean and scales variance.\\
\hlt{Key Insight}: Center at $\mu$, then subtract the fluctuation of $\overline{X}_n$.
\end{proof}
\end{theorem}

\begin{lemma} \label{lem:mean-centering-expected-square} \hlt{Expected Squared Deviation from a Sample Mean}\\
Let $X_1,\ldots,X_n$ be independent with $\mathbb{E}X_i=\mu_i$ and $\operatorname{Var}(X_i)=\sigma^2$ for $1\leq i\leq n$.\\
Set $\overline{X}\defeq \frac{1}{n}\sum\limits_{i=1}^{n}X_i$ and $\overline{\mu}\defeq \frac{1}{n}\sum\limits_{i=1}^{n}\mu_i$.
Then $\mathbb{E}(X_i-\overline{X})^2=(\mu_i-\overline{\mu})^2+\frac{n-1}{n}\sigma^2$ for $1\leq i\leq n$.
\begin{proof}
Set $U_i\defeq X_i-\overline{X}$. Then $\mathbb{E}U_i=\mu_i-\overline{\mu}$, $\operatorname{Var}(\overline{X})=\frac{\sigma^2}{n}$, and $\operatorname{Cov}(X_i,\overline{X})=\frac{\sigma^2}{n}$ by independence.\\
Hence $\operatorname{Var}(U_i)=\operatorname{Var}(X_i)+\operatorname{Var}(\overline{X})-2\operatorname{Cov}(X_i,\overline{X})=\sigma^2+\frac{\sigma^2}{n}-2\frac{\sigma^2}{n}=\frac{n-1}{n}\sigma^2$.\\
Thus $\mathbb{E}U_i^2=(\mathbb{E}U_i)^2+\operatorname{Var}(U_i)=(\mu_i-\overline{\mu})^2+\frac{n-1}{n}\sigma^2$.\\
\hlt{Structural Commentary}: $X_i-\overline{X}$ has mean $\mu_i-\overline{\mu}$ and variance $\operatorname{Var}(X_i-\overline{X})=\frac{n-1}{n}\sigma^2$.\\
\hlt{Key Insight}: $\operatorname{Cov}(X_i,\overline{X})=\frac{\sigma^2}{n}$ removes exactly one $\frac{\sigma^2}{n}$ from the raw variance $\sigma^2$.
\end{proof}
\end{lemma}

\begin{definition} \label{def:random-vector-mean-covariance-matrix} \hlt{Mean Vector and Covariance Matrix}\\
For a random vector $X=(X_1,\ldots,X_k)^t$:
\begin{enumerate}[label=(\roman*)]
\setlength{\itemsep}{0pt}
\item \hlt{Mean vector}: $\mathbb{E}X\defeq(\mathbb{E}X_1,\ldots,\mathbb{E}X_k)^t$.
\item \hlt{Covariance matrix}: $\operatorname{Var}(X)\defeq[\operatorname{Cov}(X_i,X_j)]_{i,j=1}^{k}$.
\end{enumerate}
\end{definition}

\begin{proposition} \label{prop:linear-transform-random-vector-moments} \hlt{Linear Transformations of Random Vectors}\\
Let $X\in\R^k$ satisfy $\mathbb{E}X=\mu$ and $\operatorname{Var}(X)=\Sigma$.
\begin{enumerate}[label=\roman*.]
\setlength{\itemsep}{0pt}
\item \hlt{Linear form}: for $a\in\R^k$, $\mathbb{E}(a^tX)=a^t\mu$ and $\operatorname{Var}(a^tX)=a^t\Sigma a$.
\item \hlt{Matrix transform}: for $A\in\R^{m\times k}$, $\mathbb{E}(AX)=A\mu$ and $\operatorname{Var}(AX)=A\Sigma A^t$.
\end{enumerate}
\begin{proof}
\begin{enumerate}[label=\roman*.]
\setlength{\itemsep}{0pt}
\item $\mathbb{E}(a^tX)=a^t\mathbb{E}X=a^t\mu$, and $\operatorname{Var}(a^tX)=\mathbb{E}[(a^t(X-\mu))^2]=a^t\mathbb{E}[(X-\mu)(X-\mu)^t]a=a^t\Sigma a$.
\item $\mathbb{E}(AX)=A\mathbb{E}X=A\mu$, and $\operatorname{Var}(AX)=\mathbb{E}[(AX-A\mu)(AX-A\mu)^t]=A\mathbb{E}[(X-\mu)(X-\mu)^t]A^t=A\Sigma A^t$.
\end{enumerate}
\hlt{Structural Commentary}: Linear maps push means forward and conjugate covariance matrices.\\
\hlt{Key Insight}: Center first; then move constants outside expectations.
\end{proof}
\end{proposition}

\begin{example} \label{ex:multinomial-covariance-matrix} \hlt{Multinomial Covariance Structure}\\
If $X\sim\operatorname{Multinomial}(n,p)$ with $p=(p_1,\ldots,p_k)$, then $\mathbb{E}X=np$ and $\operatorname{Var}(X_i)=np_i(1-p_i)$.\\
For $i\neq j$, $X_i+X_j\sim\operatorname{Binomial}(n,p_i+p_j)$, so $n(p_i+p_j)(1-p_i-p_j)=np_i(1-p_i)+np_j(1-p_j)+2\operatorname{Cov}(X_i,X_j)$.\\
Hence $\operatorname{Cov}(X_i,X_j)=-np_ip_j$.
\end{example}
